\chapter{数値実験}
\label{chap:experiments}

\section{実験設定}
\label{sec:setup}
%% 問題例，評価指標，比較手法，パラメータ，実行環境を明記する．
%% 調整・検証用データと最終評価用データを分け，後者で調整しない．
%% 同じテスト問題例で比較し，各手法の予算・停止条件を記す．
%% 乱数を使う手法は複数の種で実行し，種・問題例数・各問題例の実行回数を示す．
%% 問題例間のばらつきと，同じ問題例での種によるばらつきを区別する．
%% 調整の探索範囲・試行数・選択基準も記録する．

\section{実験結果}
\label{sec:results}
%% 実測値のみを載せる．平均・標準偏差の集計対象と単位を図表に明記する．

\section{考察}
\label{sec:discussion}
%% 性能，適用条件，限界を根拠とともに論じる．結果の列挙だけにしない．
